% a_evidence. Appendix file; label prefix: evi.
\section{Finite-order checks}
\noindent\genusstatus{numerically supported} The finite-order checks recorded here are computations (exact arithmetic unless stated); they support the statements at the orders and genera given, and are not proofs. The code and outputs are in \texttt{code/} of this project.

\label{sec:evi}

\subsection{Scope}
\label{sec:evi-scope}

Nothing in this appendix is a proof step: the proofs of the note are complete without it, and no statement of \cref{sec:set,sec:row,sec:lea,sec:ide,sec:reg,sec:rig} refers to it. What follows records the finite-order computations that were made while the proofs were written, as evidence that the statements are correct at the orders and weight points tested, and as a record of what was tested and what was not. The arithmetic is exact throughout, either in $\bbQ$ (rational weights, truncated series) or in a prime field $\bbF_p$ with $p=2^{127}-1$ or $p=2^{89}-1$ (random residues), unless floating point is said; a zero test in $\bbF_p$ is probabilistic. Several independent implementations were used, written separately; where a statement was checked by more than one of them this is said. A control is a deliberately wrong input (a wrong class, a missing constant, a point off $\TCY$) which must make a check fail; the controls that failed as required are listed, because they show that the checks have the power to detect an error. Weight points are written $(\alpha_0,\alpha_1,\alpha_2,\epsilon_1,\epsilon_2)$, or $(\alpha_0,\alpha_1,\alpha_2,\epsilon_1)$ with $\epsilon_2=E-\epsilon_1$ understood when the point lies on $\TCY$.

\subsection{Identities and rows}
\label{sec:evi-rows}

\emph{Identities of \cref{sec:set,sec:row,sec:lea}.} At random rational points of $\TCY$ with $\alpha_0\neq 0$ and random rational $s$, $\tau$, $z$: the amplitude identity \cref{lem:row-amp} multiplied by $(1-\tau s)^3$ (30 points); the operator identity \cref{lem:lea-op} applied to $\tau^n$, $-3\leq n\leq 6$ (12 points); and, at 25 rational points where $\delta^2$ is a rational square, for both roots of the critical polynomial: $\Ns(s)=0$; $\Lv=\wA/(\tau s)-\alpha_0$ from the superpotential; $\Lv_1+\Lv_2$, $\Lv_1\Lv_2$ and $\tau(\Lv_1-\Lv_2)=\delta$ (\cref{lem:set-eig}); $\alpha_1\alpha_2-E\Lv_2=-\wB\tau s^\ast\Lv_2$ (\cref{lem:set-Delta2}); the two forms of \cref{lem:set-N}(ii); the Hessian $\phi_\tau''(s)=-\delta/(s(1-s)U)$ and $\cW_2(U)=U^2\phi_\tau''/\tau^2$ (\cref{lem:row-hess}); the three expressions $\Delta=\epsilon_1\epsilon_2(\Lv_1-\Lv_2)(\alpha_1\alpha_2-E\Lv_1)=\alpha_1\wA\epsilon_1\epsilon_2\cW_2=-\epsilon_1\epsilon_2E^3\alg^2/\tau^2$ (\cref{prop:set-idem}, \cref{prop:lea-key}); the eikonal identity $(1-\tau)\Eik=\wB\Ns(s)$ and $\Ptr=-\Ns'(s)$ (\cref{lem:row-eik}); $\langle\chiv_1,\chiv_2\rangle=0$ and $\langle\chiv_1,\chiv_1\rangle=\Delta_1$; the three forms of $\Thf$; $\Delta_2=\epsilon_1\epsilon_2\wB\delta s^\ast\Lv_2$; the self-adjointness of $\Aop$ (\cref{prop:lea-A}) on random covectors; $\Lv_1-\dphi'=E/\tau-\wA$ (\cref{lem:row-canon}); the derivatives $\dd s^\ast/\dd\tau=s^\ast\Lv_2/\delta$ and $(\log\alg^2)_\tau=2\Rtr/\delta$ (\cref{prop:row-transport}); the identities involving $\uq$; $\eta\eta^{-1}=1$; the multiset identity of \cref{lem:set-hmult} as an identity of power sums $\sum_av_a^{1-2k}$, $k=1,2,3$; and the first-order values of \cref{rmk:row-first} ($\lrem(0)$, $(s^\ast)'(0)$, $\Delta_2'(0)=0$, the first coefficient of $\alg$). All hold. A second set of exact tests, at about 50 random rational points, covered the amplitude identity, the operator identity, $\eta^{-1}$, the identities involving $\uq$, the discriminant of $\Disc$ (\cref{lem:set-disc}(ii)), the rescaling of the rows at $\alpha_0=0$ (\cref{cor:lea-a0}) and the power-sum form of the multiset identities; a third, at 50 and 30 random points, covered the identities of \cref{lem:set-disc}(i)--(iii) ($\Disc=\cQ^2[(E-\alpha_0\tau)^2-4\alpha_1\alpha_2\tau]$, its discriminant $16\alpha_1\alpha_2\wA\wB$, $\Disc(1)=E^2$), the product identity $(\alpha_1\alpha_2-E\Lv_1)(\alpha_1\alpha_2-E\Lv_2)=\alpha_1\alpha_2\wA\wB$, and, modulo the critical polynomial $\Pq(U)$: $U\cW_U(1-U)(1-\cQ U)=\Pq(U)$, $\mathrm{disc}_U\Pq=\Disc$, the quadratic equation satisfied by $\Lv(U)$, the equality of the three forms of $\Delta$, $\Pq(1)=\wA(\cQ-1)$, $\Pq(1/\cQ)=\alpha_1(1-\cQ)/\cQ$, and $\Pq|_{\alpha_0=0}=\alpha_2(\xc\cQ U^2+\cQ(1-\xc)U-1)$ (\cref{lem:set-crit}).

\emph{The rows at the singular idempotent.} At the two points $(5/3,-4/7,13/6,9/11)$ and $(-2/5,7/4,3,1/6)$ of $\TCY$, through orders $\tau^6z^9$ and again $\tau^8z^{11}$: the point-frame row in $y=\log U$, computed from the rescaled derivatives $\wrk_k(\tau)$ through the Wick polynomials (partitions of $2n$ into $n-m$ parts $\geq 3$), multiplied by $h_{(\alpha_1,\wA,-E)}(z)$, equals $(E/(E+z))\Fhyp\Ephi\alg$ coefficientwise (\cref{thm:row-main}), and the negative powers of $z$ on the right cancel; $\rpt_{0,1}(z;0)=h_{(E,-\alpha_1,-\wA)}E/(E+z)$ (\cref{prop:row-beta}); the $z$-degree test of (L1): $\coef{t^k}(z\Gk)\prod_{j\leq k+1}(E+jz)$ has $z$-degree exactly $k$ for $1\leq k\leq 6$, resp.\ $k\leq 8$ (\cref{thm:lea-L1}, \cref{cor:lea-top}); the closed form of the leaf coefficients $\Aleaf_j$ reproduces $\coef{t^k}(z\Gk)-E\delta_{k0}$ for $k\leq 5$, resp.\ $k\leq 7$ (\cref{thm:lea-Aj}), with $\coef{t^1}(z\Gk)=(\uq/2)(1/(E+z)-1/(E+2z))$ and $\rho_1=\uq/(2E^2)$; (F2): $\rho\rr_{1,1}=E/t+\alpha_1\alpha_2/(E+z)+O(t)$ for all $z$-orders tested, with $\rr_{1,1}$ built from the $H$-row and $\Delta_1=\alpha_1\wA\epsilon_1\epsilon_2\wrk_2/\tau^2$ (\cref{cor:row-F2}, \cref{prop:lea-sigma}); row 2 of the quantum differential equation at $U_1$ (\cref{prop:lea-row2}); $\Kone=\langle\vv_1(z),\vv_1(-z)\rangle/\Delta_1=1$ (\cref{prop:lea-K1}); $\Edge_{11}(0,0)$ and $\Edge_{11}(1,0)$ have no negative powers of $t$, with values $-\epsilon_1\epsilon_2\alpha_1\alpha_2\wA\wB/(2E^2)$ and $-\epsilon_1\epsilon_2\alpha_1\alpha_2\wA\wB/(2E^3)$ at $t=0$; $\Thf(0)=0$ and $\Thf_1=-\epsilon_1\epsilon_2\alpha_1\alpha_2\wA\wB/E^2$; the flow $\dd/\dd t\,(t^2\rho^2\Edge_{11})=\Thf\,\Dg\,\Dg$ for $m+n\leq 3$ (\cref{thm:lea-flow}); the bidegree bound of (L3) for $k\leq 2$ (\cref{thm:lea-L3}); and, at the point $(0,5/2,-3/7,4/9)$, the rescaling of \cref{cor:lea-a0}: the frame-B row at $\alpha_0=0$ equals the fibrewise frame-B row at $\xc=\alpha_1/\alpha_2$ with $z\mapsto z/\alpha_2$, and $h_{(\alpha_1,\alpha_2,-E)}(z)=h_{(\cc,-\xc,-1)}(z/\alpha_2)^{-1}$. All hold at both orders.

\emph{An independent computation to higher order.} A second implementation, in $\bbF_p$ with $p=2^{127}-1$ and $p=2^{89}-1$ and truncated Laurent series in $t$ to $t^{14}$ and $z^{23}$, tested at nine weight points of $\TCY$, including points with $\alpha_1=\alpha_2$, with $\uq=0$ and with $\epsilon_1=\epsilon_2$: the formula \eqref{eq:lea-key}; (F1)--(F3); (L1) for $k\leq 13$; (L2) for $k\leq 9$; (L3) for $k\leq 9$; the flow of \cref{thm:lea-flow}; the quantum differential equation and symplecticity at both idempotents (\cref{thm:lea-qde}); the values at the regular idempotent (\cref{prop:lea-regrows}); \cref{cor:lea-a0} at $\alpha_0=0$; and the statement of \cref{rmk:lea-cy} at points off $\TCY$. All passed, and four controls failed as required: the class $(E+1,-\epsilon_1,-\epsilon_2)$ in place of $\bw$; the omission of $\hnw$; $\wA\mapsto\alpha_2$ in the rows; and $\lf_j=1/(E+j\psi)$ in place of $1/(E-j\psi)$. The same computation observed the two instances of sharpness predicted by \cref{cor:lea-top}: at $\uq=0$ the top leaf $\lf_{k+1}$ is absent from $\coef{t^k}(\rho\qt)$ for odd $k$, and at $E=9\alpha_2/7$ the coefficient of $\lf_9$ in $\coef{t^8}(\rho\qt)$ vanishes (for odd $k$ and $E=2\alpha_2$, resp.\ $E=2\alpha_1$, the product $\Psh_k$ contains the factor $(k+1)\alpha_2-\tfrac{k+1}{2}E$, resp.\ $(k+1)\wA-\tfrac{k+1}{2}E$; $\Psh_8$ contains $9\alpha_2-7E$).

\emph{Further implementations.} Three further implementations, written for the exploration that preceded the proofs, tested at five points of $\TCY$, namely
\begin{gather*}
(1,2,4,3,4),\quad(3/5,-7/3,11/2,5/7,641/210),\\
(-2,9,1/3,13/4,49/12),\quad(2/3,5/7,-9/4,3/11,-1055/924),
\end{gather*}
and the fibrewise point $(0,5/3,2/7,5/3,2/7)$; at the points
\begin{gather*}
(3,5,3,5),\ (5,11,5,11),\ (11,5,5,11),\ (3,13,5,11),\ (8,3,5,11),\\
(3,4,5,11),\ (2,7,5,11),\ (7,13,2,29),\ (17,-5,19,-8),\ (3/7,-11/5,13/9,17/4),
\end{gather*}
written as $(\alpha_1,\alpha_2,\epsilon_1,\epsilon_2)$ with $\alpha_0=E-\alpha_1-\alpha_2$ (the fifth has $\uq=0$); and at random residues: the formula \eqref{eq:lea-key} through $t^6z^{12}$ and $t^7z^{12}$; $\rpt_{0,1}(z;0)=h_{(E,-\alpha_1,-\wA)}E/(E+z)$ and $\hnw\rpt_{0,1}(z;0)=\hw E/(E+z)$ (\cref{cor:row-F2}); $\rpt_{0,2}(z;0)=h_{(\alpha_2,-E,\wB)}$ and $\vv_2(z;0)$ proportional to $(1,\alpha_1\alpha_2/E)$ with the scalar $\hreg$ (\cref{prop:lea-regrows}); the closed form of $\Aleaf_j$ for $k\leq 6$ and the exact degree $k$ of $\coef{t^k}(z\Gk)\prod(E+jz)$ for $k\leq 7$, sharp except where $\Psh_k=0$; (L3) for $k\leq 5$ (78 equations, the anti-diagonal of \cref{cor:lea-anti} being needed); (L2) for $j=1,2$, $b\leq 3$, $k\leq 5$, and the exact degrees for $j=2$, $k\leq 3$, $b\leq 2$ (\cref{prop:lea-L2}); $\Edge_{21}(b,a)=\Edge_{12}(a,b)$ (\cref{prop:lea-transp}); all $\Edge_{ij}(m,n)$ with $m+n\leq 8$ and all $\Dil_i(d)$ of $t$-order $\geq 0$ (F3); $\langle\vv_i(z),\vv_j(-z)\rangle=\delta_{ij}\Delta_i$ through $z^{13}$; row 2 of the quantum differential equation at both colours through $z^{11}$; the flow for $m+n\leq 11$; the hand values $\rho_1=\uq/(2E^2)$, the subleading term of $\Delta_1$, $\coef{t^1}(\rho\qt)=(\uq/2)(\lf_1-\lf_2)$, the value $\alpha_1\alpha_2/(E+z)$ of the $H$-row, $\Edge_{11}|_{t=0}=-(\epsilon_1\epsilon_2\alpha_1\alpha_2\wA\wB/2)\lf_1(\psi)\lf_1(\psi')$ and the rank-one $\Edge_{12}|_{t=0}$ (\cref{rmk:lea-low}, \cref{rmk:lea-legs}); the constants at $\cQ=0$ of \cref{lem:ide-norm} at both roots through $z^9$ at three points; and the $\alpha_0=0$ rescaling to the fibrewise rows. Controls: the class $(E,-\alpha_1,-\wA)$ fails at every point with $\alpha_i\neq\epsilon_i$ (it agrees only on $\TF$), the scale $E+1$ gives the full degree $13$ in the degree test, and the point frame and the class $(E,-\alpha_1,-\wA)$ both give $\Edge_{11}$ of order $t^{-2}$. At the point $(2/9,11/4,-7/6,3/5)$ the degree test of (L1) was extended: $\coef{t^k}(\rho\qt)\prod_j(E-j\psi)$ has degree exactly $k$ for $k=1,\dots,10$ through $\psi^{19}$ ($k=0$ gives the constant $E$), and the bidegree bound of (L3) holds for $k\leq 6$ and $m+n\leq 14$ with no violation in 71 to 119 tests per $k$.

\subsection{Graph regularity}
\label{sec:evi-graphs}

The frame-B graph sum \eqref{eq:set-graphsum} was evaluated per coloured stable graph, summing over decorations, in $\bbF_p$ with Laurent precision in $t$, by three implementations. The counts of stable graphs of genus $2,3,4,5$ without legs are $7,42,379,4555$; the counts of coloured stable graphs with at least one edge are $20,350,8596,284806$.

\emph{Regularity of the coloured graphs with an edge (\cref{cor:reg-main}(a)).} Genus 2: $20$ of $20$ regular at $t=0$ at the point $(1,2,4,3,4)$ and on $\TF$. Genus 3: $350$ of $350$ at four points of $\TCY$. Genus 4: $8596$ of $8596$ at three points of $\TCY$. Genus 5: $284806$ of $284806$ at the single point $(3/5,-7/3,11/2,5/7,641/210)$. A second implementation found, at three generic points of $\TCY$, that of the $22$, $352$ and $8598$ coloured stable graphs of genus $2,3,4$ exactly one is singular, the single-vertex graph at colour $1$; the first implementation was re-run at the fresh points $(7/5,-3/8,13/6,11/9)$ and $(-5/4,19/7,2/3,17/5)$ with the same result. Special points of $\TCY$ were tested to exclude an artefact of a generic sample, all with the same outcome at genus $2,3,4$: $\uq=0$ at $(1,3,2,5/2)$; $\epsilon_1=\epsilon_2$ at $(3/2,-1/3,5/7,79/84)$; $\alpha_1=\alpha_2$ at $(2/3,5/4,5/4,1/7)$; $\alpha_1=\alpha_2$ with $\epsilon_1=\epsilon_2$ at $(2,3/5,3/5,8/5)$; $3\alpha_2=E$ at $(1,3,2,7/5)$; $\wA=\alpha_1$ at $(1,3,2,1/9)$; $\TF$ at $(0,5/3,2/7,5/3)$; a point with negative weights and small $E$, $(-3,1,5/2,7)$; $\alpha_1=-\alpha_2$ at $(5,7/3,-7/3,2)$; and $\alpha_0=\alpha_1$ at $(4/3,4/3,1/5,3)$.

\emph{The single-vertex graph (\cref{cor:reg-main}(b)).} In every run above, the colour-$1$ term has the single polar term $t^{2-2g}$ with coefficient $\gapg$, and the colour-$2$ term has no pole.

\emph{Tallies.} With the pole order of a decorated graph (a coloured graph with a choice of exponents and dilaton degrees) and the bound $\sum_v(2g_v-2)_+$ of \cref{thm:reg-main}, the tallies in frame B, written as (pole order, bound) followed by the number of decorated graphs, are, identically at all generic points tested:
\begin{align*}
\text{genus }2\colon&\quad\{(0,0)\colon 20,\ (2,2)\colon 7\};\\
\text{genus }3\colon&\quad\{(0,0)\colon 504,\ (0,4)\colon 1,\ (2,2)\colon 159,\ (4,4)\colon 29\};\\
\text{genus }4\colon&\quad\{(0,0)\colon 19912,\ (0,4)\colon 25,\ (0,6)\colon 1,\ (2,2)\colon 6069,\ (4,4)\colon 1069,\ (6,6)\colon 96\}.
\end{align*}
These are the fibrewise tallies. No decorated graph exceeds its bound and none has an odd leading order. The $159$ singular decorated graphs with an edge at genus 3 (and $7138$ at genus 4) lie in graphs with a vertex of genus $\geq 2$; decorated graphs of graphs with all $g_v\leq 1$ are individually regular, while at a vertex of genus $\geq 2$ only the decoration sum at fixed colouring is regular, as the proof of \cref{thm:reg-main} predicts. The polar coefficients of individual decorated graphs depend on $\alpha$: at fixed $(\epsilon_1,\epsilon_2)=(5,11)$ and twenty values of $\alpha$, the coefficient of $t^{2-2g}$ is independent of $\alpha$, that of $t^{3-2g}$ is $\uq$ times a constant, and at genus 3 the coefficients of $t^{-2}$ and $t^{-1}$ are polynomials in $(\uq,\alpha_0^2)$ of degrees $2$ and $3$, the latter divisible by $\uq$; no proof of these observations is given.

\emph{Totals.} The total minus $\Pg\Li_{3-2g}(\cQ)/(\epsilon_1\epsilon_2)$ has zero coefficients at all non-zero powers of $t$ from $t^{1-2g}$ through $t^5$ at genus $2,3,4$ and through $t^3$ at genus 5, and its $t^0$ coefficient equals the constant-map term of \cref{rmk:rig-d0} whenever $\alpha_0\neq 0$. At a generic point $\cQ_0\in\bbF_p$ away from $\cQ=1$, with $\cQ=\cQ_0\re^t$, the total equals $\Pg\Li_{3-2g}(\cQ)/(\epsilon_1\epsilon_2)+\sF_g(0)$ through $t^4$ at genus $2,3,4$ at $(7/5,-3/8,13/6,11/9)$ and at genus 3 at $(-5/4,19/7,2/3,17/5)$; the class $(E,-\alpha_1,-\wA)$ at genus 2 and the point frame at genus 3 give the same total, as a frame change must.

\emph{Controls.} Earlier stored data were reproduced with zero mismatches: the point-frame terms on $\TCY$ at two generic points ($17$ and $378$ decorated graphs at genus 2, 3) and the frame-B terms on $\TF$ at $\xc=3/5$ and $7/3$ ($27$ and $693$). Hodge integrals computed by Mumford's recursion reproduce known values ($\int\lambda_1\lambda_2\lambda_3=1/1451520$ and $\int\lambda_2^3=1/725760$ on $\Mbar_3$, among others). Frames that must fail: the class $(E,-\alpha_1,-\wA)$ at colour $1$ gives singular edge graphs in every genus ($6$ of $6$ at genus 2, $85$ of $350$ coloured graphs or $41$ of $41$ stable graphs at genus 3, $2522$ of $8596$ at genus 4) with the total unchanged; the point frame gives $6$ of $6$ singular at genus 2 and $41$ of $41$ at genus 3, and in it every decorated graph is singular with pole order $2g-2$, the poles cancelling only in the total; the classes $\wts_\pm$ themselves at the two colours, the swapped assignment, and $\nw$ at both colours give singular edge graphs ($6$ of $6$, $41$ of $41$); the classes $(-E,\epsilon_1,\epsilon_2)$ and $(E+1,-\epsilon_1,-\epsilon_2)$, used consistently, give the correct total but singular edge graphs. The colour-$2$ class is not constrained by these tests: with the class $\bw$ at colour $1$ and the point frame, $(E,-\alpha_1,\alpha_1-E)$, $\wts_+$ or $\wts_-$ at colour $2$, all edge graphs are regular ($0$ of $6$, $0$ of $41$). Controls on the constants: omitting the constant $\hnw$, or replacing $\epsilon_1$ by $\epsilon_1+1$ in it, gives a wrong total and singular edge graphs; twisting the rows by $\hw$ while keeping the class $(E+1,-\epsilon_1,-\epsilon_2)$ gives a wrong total.

\subsection{Direct localisation}
\label{sec:evi-loc}

Independently of the R-matrix, the invariants $\GW_{g,d}$ were computed by virtual localisation on $\Mbar_g(\bbP^1,d)$ (an independent enumeration of the fixed loci with its own vertex and edge factors; Hodge integrals by Mumford's recursion), in $\bbF_p$, at points of $\TCY$ not used elsewhere, and compared with $\Pg d^{2g-3}/(\epsilon_1\epsilon_2)$: $g\leq 6$ at $d=1$ at $(7/5,-3/8,13/6,11/9)$; $g\leq 5$ at $d\leq 3$ at $(-5/4,19/7,2/3,17/5)$; $g\leq 4$ at $d\leq 4$ at $(3/2,-1/3,5/7,8/3)$; $g\leq 7$ at $d=1$ at $(3/4,-2/5,9/7,13/3)$; $g\leq 4$ at $d\leq 2$ at the point $(1,2,-3,5,-5)$ with $E=0$ and at the point $(2/3,5/4,5/4,1/7)$ with $\alpha_1=\alpha_2$. All agree, including genus $0$ and genus $1$ (expected $1/(d^3\epsilon_1\epsilon_2)$ and $\Ph{1}/(d\epsilon_1\epsilon_2)$, \cref{prop:rig-g0}, \cref{lem:rig-g1}). The control off $\TCY$ at $(7/5,-3/8,13/6,11/9,2)$ fails for every $g\geq 1$. A separate direct-localisation computation at $d=1$, $g\leq 2$, confirmed $\epsilon_1\epsilon_2\GW^{\Tmax}_{1,1}$ of \cref{lem:rig-g1} on the five-torus and $\epsilon_1\epsilon_2\GW_{2,1}=\Ph{2}$ on $\TCY$; both agree with the genus $\leq 2$ part of the five-torus expansion displayed in \cite{BS} off $\TCY$. \draftnote{BS label: the displayed five-torus invariants of degree one in the resolved-conifold section, eq: res con GW 2.}

\subsection{Rational reconstruction in \texorpdfstring{$\cQ$}{Q}}
\label{sec:evi-rational}

From $60$ to $70$ evaluations at points $\cQ\in\bbF_p$, the frame-B graph terms were reconstructed as rational functions of $\cQ$ at a generic point of $\TCY$. Genus 2: the total is $N/(1-\cQ)^2$ with $\deg N=2$; each of the $6$ edge graphs is $N/\Disc^4$ with $\deg N=8$; the vertex graph is $N/(\Disc^4(1-\cQ)^2)$ with $\deg N=10$. Genus 3: the total is $N/(1-\cQ)^4$ with $\deg N=4$; each of the $41$ edge graphs is $N/\Disc^8$ with $\deg N=16$; the vertex graph is $N/(\Disc^8(1-\cQ)^4)$ with $\deg N=20$. So the individual graphs have poles on $\Disc=0$, of orders $4$ and $8$, which cancel in the total (\cref{lem:rig-disc}); the edge graphs are regular at $\cQ=1$ and all graphs are regular at $\cQ=\infty$ (\cref{lem:rig-inf}); and the total has the pole structure of $\Pg\Li_{3-2g}/(\epsilon_1\epsilon_2)+\sF_g(0)$ (\cref{thm:rig-main}). In the point frame at genus 2 every graph has a pole of order $2$ at $\cQ=1$.

\subsection{The rigidity argument}
\label{sec:evi-rig}

Exact rational arithmetic at random rational weights:
\begin{enumerate}
\item[(A)] $\Ssh(y)=y/(2\sinh(y/2))=1-y^2/24+7y^4/5760-31y^6/967680+127y^8/154828800-\dots$; $\Ph{2}$ and $\Ph{3}$ as polynomials; $(2g-3)!\Pg/(\epsilon_1\epsilon_2)=\beta_2,\beta_4$ as displayed in \cite{BS} (\cref{lem:reg-beta}); $\gapg$ at $(\epsilon_1,\epsilon_2)=(\xc,1)$, $g=2$, equals $(7+10\xc^2+7\xc^4)/(5760\xc)$. \draftnote{BS label: eq:betapol.}
\item[(B)] At 50 points: $\Disc=\cQ^2[(E-\alpha_0\tau)^2-4\alpha_1\alpha_2\tau]$; $\mathrm{disc}_{\cQ}\Disc=16\alpha_1\alpha_2\wA\wB$; $\Disc(1)=E^2$; $(\alpha_1\alpha_2-E\Lv_1)(\alpha_1\alpha_2-E\Lv_2)=\alpha_1\alpha_2\wA\wB$ (\cref{lem:set-disc}, \cref{prop:set-idem}).
\item[(C)] In floating point at $\cQ=10^6$ and $10^9$: $U_1\to(\alpha_1-\alpha_2)/\wB$, $\cQ U_2\to-\alpha_2/(\alpha_1-\alpha_2)$, $\Lv_i\to\alpha_1,\alpha_2$, $\Delta_1\to-\alpha_1\wB(\alpha_1-\alpha_2)\epsilon_1\epsilon_2$, $\Delta_2\to\alpha_2\wA(\alpha_1-\alpha_2)\epsilon_1\epsilon_2$ (\cref{lem:rig-inf}, \eqref{eq:rig-Dinf}).
\item[(D)] The simple-root lifting of \cref{lem:rig-root} iterated in $\Kred\fps{\alpha_0}$ to $O(\alpha_0^8)$ at five random $(\alpha_1,\alpha_2)$: $\Disc(\cQ_1)=0$ to that order, with the displayed coefficients of $\alpha_0^2,\alpha_0^3,\alpha_0^4$.
\item[(E)] $\Li_{-n}=\Aeul_n/(1-\cQ)^{n+1}$ with $\Aeul_n(0)=0$ and $\deg\Aeul_n=n$ ($\Aeul_1=\cQ$, $\Aeul_2=\cQ+\cQ^2$, $\Aeul_3=\cQ+4\cQ^2+\cQ^3$, $\Aeul_4=\cQ+11\cQ^2+11\cQ^3+\cQ^4$); $\Li_{-n}(1/\cQ)=(-1)^{n+1}\Li_{-n}(\cQ)$ for $n=1,\dots,7$ at ten random $\cQ$ each; agreement with the series at $\cQ=1/3$ to $10^{-30}$; $\Li_0(1/\cQ)=-1-\Li_0(\cQ)$; the expansion of $\Li_{-n}(\re^t)$ at $t=0$ has polar part exactly $(-1)^{n+1}n!\,t^{-n-1}$ for $n=1,\dots,6$ and no other negative power, with $\Li_{-1}(\re^t)=t^{-2}-1/12+O(t^2)$ and $\Li_{-3}(\re^t)=6t^{-4}+1/120+O(t^2)$ (\cref{lem:rig-li}).
\item[(F)] The multiset identities as power sums $\sum_av_a^{1-2k}$, $k=1,2,3$, at 20 points: $\nw+(E,-\alpha_1,-\wA)=(E,-\epsilon_1,-\epsilon_2)$ (\cref{lem:set-hmult}); at $\alpha_0=0$, $\nw-\bw=(\alpha_1,\alpha_2,-E)$ (\cref{lem:rig-3Tid}); at $\cQ=\infty$, $(\alpha_2-\alpha_1,-\wA,\wB)+\nw=\wts'_-$ and $(\alpha_2,-\alpha_1,\alpha_1-\alpha_2)+\nw=\wts'_+$ (\cref{prop:rig-flop}).
\end{enumerate}
The $\alpha_0$-adic step of \cref{lem:rig-disc} and the regularity at $\cQ=\infty$ of \cref{lem:rig-inf} were tested only through the totals and the rational reconstruction of \cref{sec:evi-rational}.

\subsection{The flop and the constant term}
\label{sec:evi-flop}

The identity $\sF_g(0)=\sF^{\Zfl}_g(0)$ of \cref{prop:rig-flop}(iii), with both sides computed as the constant-map Hodge integrals of \cref{rmk:rig-d0} at the weights $\wts_\pm$ and $\wts'_\pm$, holds for $g=2,\dots,6$ at two points of $\TCY$ and fails for $g=2,\dots,6$ at the point $(1,2,4,3,5)$ off $\TCY$. The critical points at $\sQ=0$, the limits of $\cW_2(U_i)$ and the weights $\wts'_\pm$ were also checked exactly by a second computation. The constant term $\sF_g(0)$ itself depends on $\alpha$: at fixed $(\epsilon_1,\epsilon_2)$ and three values of $\alpha$ it does not take the same value at the three for $g=2,3,4$ (\cref{rmk:rig-F0}).

\subsection{Off the Calabi--Yau torus}
\label{sec:evi-offCY}

With the class $(E_\alpha,-\epsilon_1,-\epsilon_2)$, $E_\alpha\coloneqq\alpha_0+\alpha_1+\alpha_2\neq\epsilon_1+\epsilon_2$, at the points $(2,3,4,5,11)$, $(3,3,4,5,11)$, $(1,7,2,5,3)$ and $(1,2,4,3,5)$: the singular row still freezes, $\rr_{0,1}(z;0)=E_\alpha/(E_\alpha+z)$, and $\Delta_1t^2\to-E_\alpha^3\epsilon_1\epsilon_2$ (\cref{rmk:lea-cy}); at genus 2 the edge graphs are regular but the vertex graph acquires a $t^{-1}$ term; at genus 3 four stable graphs are singular (those with a vertex of genus $\geq 2$; $4$ of $350$ coloured graphs), at genus 4 fifty-four of $8596$ coloured graphs are singular, the vertex graph has all polar orders, and the total is not $\Pg\Li_{3-2g}/(\epsilon_1\epsilon_2)+\sF_g(0)$. Direct localisation off $\TCY$ fails for every $g\geq 1$ (\cref{sec:evi-loc}), and the flop identity fails (\cref{sec:evi-flop}). The leaf structure survives off $\TCY$ in these tests; what is lost is the vanishing theorem, which needs the sum-zero class $\bw$ (\cref{thm:reg-L}, \cref{rmk:lea-cy}).

\subsection{What the checks do not cover}
\label{sec:evi-limits}

All checks are at finite genus ($g\leq 5$ for the graph sums, with a single point at genus 5; $g\leq 7$ for direct localisation at $d=1$; $g\leq 6$ for the flop identity), at finite orders in $t$, $z$, $\psi$, and at finitely many weight points; the zero tests in $\bbF_p$ are probabilistic. The exact agreement in $\cQ$ on $\TCY$ exists only through the rational reconstruction at genus $2,3$. Not tested at all: the stability argument (Fl6) and the descent of \cref{prop:ide-teleman}, the analytic input \cref{lem:row-betaphase} beyond the values it predicts, and the identification of \cref{prop:rig-flop}(iii) beyond the Hodge-integral identity itself. The two implementations of the graph sum agree with each other on the points compared, and with the direct localisation, which does not use the R-matrix at all.
